\documentclass[12pt]{article}
\usepackage[utf8]{inputenc}
\usepackage[T1]{fontenc}
\usepackage{hyperref}
\title{Subtiaba Language Resources for K-12 Curriculum Development}
\author{Compiled by Claude Sonnet 5.5 (claude-sonnet-5-5) for Andr\'as Kornai}
\date{}
\begin{document}
\maketitle

\section{Introduction}
Subtiaba (ISO 639-3: sut; Glottocode: subt1250) is a dormant Otomanguean language of the Pacific slope of Nicaragua, historically spoken in the Subtiaba district of Le\'on. It is grouped with Tlapanec in the Supanecan branch. Sources report it was extinct by the 1920s. My search found no modern published grammar, dictionary, primer or school programme. What survives is mainly 19th-century manuscript material, so any curriculum would have to be a revitalization or heritage project built on archival sources.

\section{Primary Educational Resources}
No primers, textbooks or school materials for Subtiaba were found in the search results. The language is not known to be taught in schools.

\section{Linguistic and Reference Materials}
\begin{itemize}
\item Subtiaba language (Wikipedia overview, with classification and references): \url{https://en.wikipedia.org/wiki/Subtiaba_language}
\item C. Hermann Berendt, 1874 manuscript on Subtiaba vocabulary, pronunciation and grammar (University of Pennsylvania), including notes copied from a grammar attributed to Francisco D\'iaz Zapata: \url{https://openn.library.upenn.edu/Data/0002/html/mscoll700_item242.html}
\item Berendt's comparative vocabulary (Ms. Coll. 700, item 144): about 750 words, 42 numerals and 81 phrases in Spanish with Subtiaba, Chiapanec and Mangue equivalents: \url{https://openn.library.upenn.edu/Data/0002/html/mscoll700_item144.html}
\item Colenda Digital Repository record (University of Pennsylvania): \url{https://colenda.library.upenn.edu/catalog/81431-p39c6sn0w}
\item Rosetta Panglossia entry for sut: \url{https://rosettapanglossia.longnow.org/index.php?title=Sut}
\end{itemize}

\section{Religious and Community Resources}
None found in the search results.

\section{Digital and Audio Resources}
No audio, video or software resources were found.

\section{Curriculum Development Context}
There are no known native speakers, no modern description and no teaching materials. Any curriculum would have to start from the archival manuscripts above, ideally with linguists and descendant community members in Le\'on, and would be a reconstruction effort. Content in the language itself should be kept to what the sources actually attest. Details of the manuscripts should be checked against the repositories directly. I did not open the manuscript images.

\end{document}